% Zone „Das kennst du schon“ – aus bank/rekonstruktion-von-bestaenden/zone.jsonl (zusammenbau v0.1)
\einheitenkopf[zone][Kennst du schon]{Das kennst du schon}
\setcounter{aufgabe}{0}

%% TODO Titel: Ich-kann-Satz fehlt in der Bank (Platzhalter: Kettenname) – Nr. 1
%% TODO Anweisung: ein Satz über der ersten Teilaufgabe fehlt in der Bank – Nr. 1
\begin{aufgabe}{Bestimmte Integrale berechnen (Stammfunktion, Hauptsatz, auch mit Rechner)}
\begin{teile}
\teil $∫_0^3 2t\,\mathrm{d}t$? \leerfeld
\teil $∫_1^3 (1{,}5t^2 - 2t)\,\mathrm{d}t$? \leerfeld
\end{teile}
\end{aufgabe}

%% TODO Titel: Ich-kann-Satz fehlt in der Bank (Platzhalter: Kettenname) – Nr. 2
%% TODO Anweisung: ein Satz über der ersten Teilaufgabe fehlt in der Bank – Nr. 2
\begin{aufgabe}{Die Rate als Ableitung lesen}
\begin{teile}
\teil $V(t)$ ist das Wasservolumen in Litern nach $t$ Minuten, $V'(3) = 4$. Was bedeutet das?
\teil $B(t)$ ist die Zahl der Gäste in einem Bad nach $t$ Stunden, $B'(2) = -150$. Was bedeutet das?
\end{teile}
\end{aufgabe}

%% TODO Titel: Ich-kann-Satz fehlt in der Bank (Platzhalter: Kettenname) – Nr. 3
%% TODO Anweisung: ein Satz über der ersten Teilaufgabe fehlt in der Bank – Nr. 3
\begin{aufgabe}{Vorzeichen einer Funktion über Nullstellen bestimmen}
\begin{teile}
\teil Wo ist $f(t) = t - 3$ negativ, wo positiv?
\teil Wo ist $f(t) = -0{,}5 \cdot (t - 2) \cdot (t - 6)$ positiv?
\end{teile}
\end{aufgabe}

%% TODO Titel: Ich-kann-Satz fehlt in der Bank (Platzhalter: Kettenname) – Nr. 4
%% TODO Anweisung: ein Satz über der ersten Teilaufgabe fehlt in der Bank – Nr. 4
\begin{aufgabe}{Einheiten umrechnen}
\begin{teile}
\teil 3\,m³ in Liter? \leerfeld[l]
\teil 72\,km/h in m/s? \leerfeld m/s
\end{teile}
\end{aufgabe}

%% TODO Titel: Ich-kann-Satz fehlt in der Bank (Platzhalter: Kettenname) – Nr. 5
%% TODO Anweisung: ein Satz über der ersten Teilaufgabe fehlt in der Bank – Nr. 5
\begin{aufgabe}{Rate mal Zeit als Menge deuten (je-desto, proportionale Grundvorstellung)}
\begin{teile}
\teil Ein Hahn liefert 6 Liter je Minute. Wie viel Wasser in 5 Minuten? \leerfeld[l]
\teil Erst fließen 12 Minuten lang 2{,}5 Liter je Minute zu, dann 8 Minuten lang 1{,}5 Liter je Minute. Wie viel insgesamt? \leerfeld[l]
\end{teile}
\end{aufgabe}

%% TODO Titel: Ich-kann-Satz fehlt in der Bank (Platzhalter: Kettenname) – Nr. 6
%% TODO Anweisung: ein Satz über der ersten Teilaufgabe fehlt in der Bank – Nr. 6
\begin{aufgabe}{Einheiten umrechnen – Fehler finden}
\begin{teile}
\teil Ein Zug fährt mit 240\,km/h. Tim rechnet für drei Minuten: \rechnung{240 \cdot 3 &= 720\,\mathrm{km}} Finde den Fehler und rechne richtig.
\end{teile}
\end{aufgabe}

%% TODO Titel: Ich-kann-Satz fehlt in der Bank (Platzhalter: Kettenname) – Nr. 7
%% TODO Anweisung: ein Satz über der ersten Teilaufgabe fehlt in der Bank – Nr. 7
\begin{aufgabe}{Einheiten umrechnen – selbst rechnen}
\begin{teile}
\teil Ein Auto fährt mit 90\,km/h. Welche Strecke legt es in 20 Minuten zurück? \leerfeld[km]
\end{teile}
\end{aufgabe}
